% problem-000219 12 金属卡宾热裂解
\section*{12. 金属卡宾热裂解}
\textit{下列为分问。}

\subsection*{12-1}
化合物A真空热裂解，经卡宾中间体B转化为产物C。画出中间体B的结构。

（图：）

\subsection*{12-2}
⾦属络合物可以在温和条件下促进炔烃的转化反应。在[RuCl_{2}(C_{6}H_{6})PPh_{3}]或 CpRu(PPh_{3})Cl 的存在下，炔烃经由⾦属卡宾中间体D发⽣⽔合反应⽣成醛F（以[Ru]表⽰催化剂）。画出中间体D、E和最终产物F的结构。（若有形式电荷请标明，不考虑⽴体化学）

（图：）

\subsection*{12-3}
根据以上信息，画出下图中最终产物G的结构。

（图：）

\subsection*{12-4}
画出以下反应过程中的中间体I和J的结构。（若有形式电荷请标明，不考虑⽴体化学）

（图：）

\subsection*{12-5}
画出以下反应过程中的中间体L和M的结构。（若有形式电荷请标明，不考虑⽴体化学）

（图：）

\subsection*{12-6}
画出反应中间体0和P的结构（若有形式电荷请标明，不考虑立体化学），并解释以下反应生成不同环系的原因。

（图：）
